\chapter{Metody}

Tato kapitola popisuje metody použité při zpracování práce. Analytická část se opírá o~rešerši, analýzu datových zdrojů a~komparaci, realizační část o~metodiku vývoje aplikace. Rešerše probíhala od října 2024 do dubna 2026, analýza datových zdrojů a~komparace v~dubnu 2026.

\section{Rešerše}

Rešerše pokrývala čtyři tematické okruhy: kvalitu vnitřního prostředí, otevřená data, open source a~existující řešení pro sběr dat (aplikace a~zařízení).

\FloatBarrier
\subsection{Kvalita vnitřního prostředí}

Rešerše směřovala od identifikace základních zdrojů o~IEQ/IAQ k~vyhledání definic jednotlivých složek a~příslušných standardů a~norem.

Základní soubor zdrojů byl identifikován nástrojem \href{https://www.proquest.com/}{ProQuest}. Použita byla anglická\footnote{Veškeré překlady z~anglického jazyka byly provedeny buď za pomocí jazykové výbavy autora nebo s~pomocí online překladačů Google Translate a~DeepL a~v~krajních případech za pomocí LLM. V~případě, že překlad do češtiny byl nejednoznačný nebo nepřesný, byla ponechána originální verze v~angličtině. Ponechání v~angličtině bylo provedeno i~v~případě, že se jednalo o~běžné odborné termíny v~dané oblasti.} klíčová slova \textit{Indoor Environment Quality}, \textit{IEQ}, \textit{Indoor Air Quality}, \textit{IAQ}, \textit{productivity} a~\textit{health} v~kombinaci s~operátory AND a~OR, řazená podle relevance a~omezená na posledních pět let.

Standardy, normy a~doporučení byly vyhledány přes Google Scholar a~právní informační systém \href{https://www.aspi.cz/menu/search}{ASPI}, doplňkově byla využita databáze \href{https://www.ieqguidelines.org/}{IEQ Guidelines} shromažďující požadavky z~celého světa. Klíčová slova v~češtině i~angličtině pokrývala vnitřní prostředí, kvalitu vzduchu, větrání, hygienické požadavky, tepelný komfort, osvětlení, denní osvětlení a~hluk (např. \textit{WHO guidelines indoor air quality}, \textit{ventilation buildings}, \textit{environmental noise health}).

\FloatBarrier
\subsection{Otevřená data}

Cílem bylo vymezit otevřená data v~české legislativě a~zjistit způsob jejich publikace. V~systému \href{https://www.aspi.cz/menu/search}{ASPI} byl termínem \textit{otevřená data} (řazeno dle relevance, omezeno na platné a~účinné předpisy) identifikován zákon 106/1999 Sb., o~svobodném přístupu k~informacím, který zavádí pojmy \textit{otevřená data} a~\textit{otevřená formální norma} (OFN). Z~referencí v~zákoně byly dohledány vazby na legislativu EU.

Vyhledávačem \href{https://duckduckgo.com/}{DuckDuckGo} byly podle termínů z~legislativy identifikovány Národní katalog otevřených dat \href{https://data.gov.cz/}{data.gov.cz} a~evropský portál \href{https://data.europa.eu/}{data.europa.eu}, který posloužil jako klíčový zdroj informací o~licencích. Relevance Národního katalogu a~OFN pro soukromé subjekty byla ověřena e-mailovou konzultací s~koordinátorem Národního katalogu z~Digitální a~informační agentury (korespondence viz příloha~\ref{app:korespondence-dia}).

\FloatBarrier
\subsection{Open source}

Pro vymezení pojmu open source byly využity \href{https://scholar.google.cz/}{Google Scholar} a~\href{https://duckduckgo.com/}{DuckDuckGo} s~termíny \textit{open source}, \textit{definition} a~\textit{open source definition}. Tím byl identifikován dokument \emph{The Open Source Definition} organizace \href{https://opensource.org/}{Open Source Initiative} (OSI), uznávané autority v~oblasti, který stanovuje kritéria pro open source licence (mimo jiné volnou redistribuci a~přístup ke zdrojovému kódu). Přes Google Scholar byly doplněny akademické zdroje o~kvalitě open source projektů.

\FloatBarrier
\subsection{Existující aplikace pro sběr IEQ dat}

Existující aplikace byly identifikovány funkcí \emph{Hluboký výzkum} nástroje ChatGPT, která umožňuje vícestupňové prohledávání internetu. Použit byl prompt: \textit{Find applications that collect Indoor Environment Quality (IEQ) data and publish the data as open data for everyone to use. Focus on finding applications that are open source (if you are not sure, it does not matter, list it anyways, I will confirm myself).} Takto bylo nalezeno pět aplikací, jejichž relevance byla následně ručně ověřena.

\FloatBarrier
\subsection{Zařízení pro sběr dat}

Přehled senzorových modulů a~hotových zařízení pro monitorování IEQ vznikl primárně z~výsledků analýzy existujících aplikací (kapitola~\ref{ch:aplikace-analyza}). Základ tvořil seznam podporovaných senzorů projektu sensor.community \parencite{SensorCommunityDevices}, doplněný o~senzory a~zařízení, pro něž poskytují firmware, návody nebo přímou integraci platformy OpenAQ, openSenseMap a~IndoorCO2Map. Doplňkově byly využity zoterové kolekce \textit{senzory} a~\textit{Podporovane sensor.community} s~technickými listy dalších senzorů. Průmyslová zařízení Elvaco Sense 100W a~300W byla zařazena na základě přímé znalosti autora z~předchozího zaměstnání, kde se podílel na aplikaci tyto senzory využívající.

Pro každý senzor nebo zařízení byl dohledán technický list nebo produktová stránka výrobce a~z~nich zjištěny měřené veličiny IEQ, komunikační rozhraní, přesnost a~orientační cena. Zdroje byly zařazeny do bibliografické databáze a~použity pro přehledové tabulky.

\section{Analýza datových zdrojů}

Při identifikaci měřených veličin projektu sensor.community byly analyzovány datové archivy ke stažení z~\href{https://archive.sensor.community/}{archivu} a~rozhraní API. Protože dokumentace projektu neuvádí úplný a~jednoznačný seznam veličin, byly jejich unikátní názvy získány automatizovanou extrakcí z~dat za březen 2026.

\FloatBarrier
\subsection{Validace souladu datových formátů s~otevřenými formálními normami}
\label{subsec:validace-ofn}

Cílem bylo ověřit, do jaké míry datové archivy a~API analyzovaných aplikací odpovídají doporučením OFN z~katalogu špatné praxe portálu data.gov.cz \parencite{PrikladySpatnePraxe2022}. Validaci podle detailního zadání autora provedl asistent s~umělou inteligencí (Claude) s~využitím nástrojů příkazové řádky (\texttt{curl}, \texttt{jq}, \texttt{file}, \texttt{python3}).

Postup měl tři fáze. Nejprve byl z~obsahu deseti podstránek katalogu špatné praxe sestaven soubor hodnoticích kritérií (právní překážky, dokumentace a~metadata datových sad, zabezpečení přenosu, chyby formátů CSV a~JSON a~reprezentace prostorových dat). Poté byly staženy a~ověřeny vzorky datových archivů všech čtyř aplikací (formát a~kódování, struktura, validita JSON, u~IndoorCO2Map i~celistvost dat). Nakonec byly u~API zkontrolovány HTTP hlavičky, způsob autentizace a~reprezentace numerických hodnot, časových razítek a~prostorových souřadnic ve vzorové odpovědi. Zjištění byla porovnána se sestavenými kritérii a~zapracována do sekce~\ref{subsec:soulad-ofn}.

\section{Komparace}

Komparace shrnuje výsledky předchozích analýz do přehledových srovnání ve třech oblastech: zákonných limitů pro IEQ, aplikací pro sběr senzorických dat a~zařízení pro sběr dat.

\FloatBarrier
\subsection{Komparace zákonných limitů pro IEQ}

Limitní a~doporučené hodnoty pro jednotlivé složky IEQ získané rešerší byly porovnány napříč zdroji -- českou legislativou, předpisy EU a~doporučeními WHO -- s~cílem zjistit shody a~rozdíly v~přípustných hodnotách a~rozsah pokrytí jednotlivých veličin. Výsledné přehledové tabulky jsou součástí přílohy~\ref{app:legislativa} a~slouží jako podklad pro formulaci požadavků na měřené veličiny vlastní aplikace.

\FloatBarrier
\subsection{Komparace aplikací pro sběr senzorických dat}

Cílem komparace bylo syntetizovat výsledky individuálních analýz čtyř aplikací (Sensor.Community, OpenAQ, openSenseMap a~IndoorCO2Map) a~identifikovat vzory, odlišnosti a~mezery jako vstup pro analýzu požadavků vlastní aplikace.

Komparace proběhla v~pěti dimenzích. První dimenzí bylo pokrytí veličin relevantních pro IEQ, systematizované do tabulky~\ref{tab:aplikace-srovnani-data} a~umožňující přímé porovnání podporovaných veličin napříč kategoriemi IEQ. Druhou dimenzí byly technické parametry (komunikační model, formát API a~archivu, autentizace, přístup k~definici veličin a~závislost na hardwaru), zachycené v~tabulce~\ref{tab:aplikace-srovnani-technicke}. Třetí dimenzí byla webová rozhraní hodnocená na základě vlastního pozorování (vizualizace dat, filtry a~analytické funkce), čtvrtou přístup k~hardwarové podpoře z~hlediska škálovatelnosti a~konzistence dat. Pátou dimenzí, přesahující rámec standardní komparace, byl soulad s~otevřenými formálními normami, hodnocený empirickou validací podle sekce~\ref{subsec:validace-ofn}. Výsledky jsou shrnuty v~sekci~\ref{sec:srovnani-aplikaci}.

\FloatBarrier
\subsection{Komparace zařízení pro sběr dat}

Identifikovaná zařízení byla rozdělena na senzorové moduly a~hotová zařízení, protože se obě skupiny zásadně liší způsobem nasazení, cenou i~mírou otevřenosti. Hotová zařízení byla dále členěna na spotřebitelský a~průmyslový segment vykazující odlišné komunikační protokoly a~přístupy k~otevřenosti dat. Pro každou kategorii byla sestavena přehledová tabulka s~kritérii výrobce, měřené veličiny IEQ, komunikační rozhraní, deklarovaná přesnost a~orientační cena, u~hotových zařízení navíc s~hodnocením přístupu k~otevřenému API. Výsledné tabulky jsou součástí kapitoly~\ref{sec:zarizeni} a~slouží jako podklad pro požadavky na vlastní aplikaci.

\FloatBarrier
\section{Metodika vývoje aplikace}
\label{sec:metodika-vyvoje}

Realizace aplikace stavěla na třech vzájemně provázaných přístupech -- návrhu řízeném doménou (Domain-Driven Design, DDD), vývoji řízeném testy (Test-Driven Development, TDD) a~agentickém programování. Tato sekce shrnuje jejich roli jako metody, konkrétní projevy v~kódu popisuje kapitola~\ref{ch:implementace} a~vyhodnocení testů kapitola~\ref{ch:testovani}.

\textbf{Návrh řízený doménou} (DDD) určil hrubé členění systému. Problémová oblast byla rozdělena na ohraničené domény (\emph{bounded contexts}), z~nichž se každá stala samostatnou službou s~vlastním datovým modelem a~úložištěm. Jednotná terminologie domény (uživatel, budova, místnost, senzor, veličina, měření) se promítá do názvů v~kódu i~v~datovém modelu, takže slovník analýzy, návrhu a~implementace zůstává jednotný.

\textbf{Vývoj řízený testy} (TDD) určil postup psaní kódu. Každá jednotka funkcionality vznikala v~cyklu \uv{nejprve test, poté implementace, nakonec úprava} (\emph{red--green--refactor}). Test psaný před implementací vynucuje promyšlení hraničních scénářů předem a~poskytuje ověřitelné kritérium správnosti.

\textbf{Agentické programování} využilo asistenta s~umělou inteligencí (Claude) ke generování testů i~implementačního kódu na základě zadání autora. Postup TDD přitom zůstal zachován -- nejprve vznikaly testy vymezující očekávané chování, teprve poté kód, který jim měl vyhovět. Úloha autora spočívala v~zadávání požadavků a~v~rozhodování o~přijetí, úpravě nebo zamítnutí každého navrženého řešení, za jehož výsledek nesl odpovědnost.
